(1) 中心力为 $F_r=-\kappa/r^2+2\beta/r^3$，角动量守恒为 $L=mr^2\dot\phi$。令 $u=1/r$ 后 Binet 方程 $u''+u=-mF_r/(L^2u^2)$，即 $u''+[1+2m\beta/L^2]u=m\kappa/L^2$。这里撇表示对 $\phi$ 的导数，附加势改变的是角向振荡频率，而非把 $\kappa$ 简单重整化。
(2) 令 $\nu=\sqrt{1+2m\beta/L^2}>0$，积分得 $u(\phi)=m\kappa/(L^2\nu^2)[1+e\cos\nu(\phi-\phi_0)]$。因此 $r(\phi)=p/[1+e\cos\nu(\phi-\phi_0)]$，$p=(L^2+2m\beta)/(m\kappa)$；$0\le e<1$ 为有界非碰撞轨道，$e=0$ 给出圆轨道 $r_c=p$。轨道径向最小值发生在余弦为一处。
(3) 连续两个近心点满足相位 $\nu\phi$ 增加 $2\pi$，故 $\Delta\phi=2\pi/\nu$，相对开普勒情形的近心进动为 $\delta\phi=2\pi(\nu^{-1}-1)$。$\beta>0$ 时 $\nu>1$ 因而退行；$-L^2/(2m)<\beta<0$ 时 $\nu<1$ 因而顺行。若 $\nu$ 为有理数才可能在有限周期后形成闭合轨道，普通参数不满足此条件。
(4) 径向有效势为 $U_{\rm eff}=(L^2+2m\beta)/(2mr^2)-\kappa/r$；有正的离心系数时最小值为 $-\frac{m\kappa^2}{2(L^2+2m\beta)}$，有界运动要求 $U_{\min}\le E<0$，并有 $e^2=1+2E(L^2+2m\beta)/(m\kappa^2)$。若 $L^2+2m\beta\le0$，短程势不再产生阻止坠心的正势垒，稳定圆轨道消失。
